\input{../tex-inputs/latex-leseansicht-vorspann}

\section[Gustav Schwarzkopf an Arthur Schnitzler, {[}3. 10. 1902?{]}]{L04241 Gustav Schwarzkopf an Arthur Schnitzler, {[}3. 10. 1902?{]}}
\nopagebreak\mylabel{L04241v}
\rehead{ }\normalsize\beginnumbering\briefempfaengerindex{Schnitzler, Arthur@\textsc{Schnitzler, Arthur}!zzzSchwarzkopf, Gustav@\emph{von Gustav Schwarzkopf}!1902-10-032@{{[}3. 10. 1902?{]}}|(be}
\toendnotes[C]{\smallbreak\pagebreak[2]}
\correspDesc{Versand  durch Gustav Schwarzkopf am [3. 10. 1902?] in Wien
\newline{}Erhalt  durch Arthur Schnitzler im Zeitraum [3. 10. 1902 – 6. 10. 1902?] in Wien}\toendnotes[C]{\smallbreak}
\Standort{CUL, Schnitzler, B 96a.}
\physDesc{Visitenkarte, 284 Zeichen
\newline{}Handschrift: schwarze Tinte, deutsche Kurrent}\toendnotes[C]{\smallbreak}
\pstart
           \noindent{}{\pb}Lieber Arthur, es iſt
               mir leider \label{K_L04241-1v}\edtext{nicht möglich zum
                  Nachtmahl}{\lemma{\textnormal{\emph{nicht … Nachtmahl}}}\Cendnote{\textnormal{Zur Datierung siehe XXXX Auszeichnungsfehler: Dokument L04178 nicht gefunden.}}}\label{K_L04241-1} zu ko{\geminationm}en; ich muß in’s Volkstheater\oindex{Wien@\textbf{Wien}!VII., Neubau@\textbf{VII., Neubau}!Volkstheater@\textbf{Volkstheater}, \emph{Theater}|pw} zur »Freundin\pwindex{Brociner, Marco 20.\,10.\,1852 Iași – 12.\,4.\,1942 Wien@\textsc{Brociner, Marco} (20.\,10.\,1852 Iași – 12.\,4.\,1942 Wien), \emph{Schriftsteller, Journalist, Kritiker}!Freundin@\strich\emph{Die Freundin}|pw}\eventindex{Volkstheater@\textbf{Volkstheater}!Uraufführung von Die Freundin, 20.9.1902@Uraufführung von Die Freundin, 20.9.1902|pwv}«. Brociner\pwindex{Brociner, Marco 20.\,10.\,1852 Iași – 12.\,4.\,1942 Wien@\textsc{Brociner, Marco} (20.\,10.\,1852 Iași – 12.\,4.\,1942 Wien), \emph{Schriftsteller, Journalist, Kritiker}|pw} hat mir heute höchſt
               »perſönlich« Karten gebracht. Da ka{\geminationn} ich doch nicht
               anders –\pend
           
\pstart
           Bitte entſchuldigen Sie mich und empfehlen Sie mich beſtens Ihrer Frau Mama\pwindex{Schnitzler, Louise 8.\,7.\,1840 Kőszeg – 9.\,9.\,1911 Wien@\textsc{Schnitzler, Louise} (8.\,7.\,1840 Kőszeg – 9.\,9.\,1911 Wien)|pwv}.\pend
           \pstart Herzlichſt Ihr\pend{}
\pstart
           \centering{}{\pb}\textcolor{gray}{\textbf{Gustav Schwarzkopf}}\pend
           \selectlanguage{ngerman}\endnumbering\briefempfaengerindex{Schnitzler, Arthur@\textsc{Schnitzler, Arthur}!zzzSchwarzkopf, Gustav@\emph{von Gustav Schwarzkopf}!1902-10-032@{{[}3. 10. 1902?{]}}|)be}\mylabel{L04241h}
\begin{anhang}
\end{anhang}\newcommand{\dateiname}{L04241}\newcommand{\titel}{Gustav Schwarzkopf an Arthur Schnitzler, [3. 10. 1902?]}\newcommand{\editorInnen}{Herausgegeben von Jahnke, SelmaMüller, Martin Anton}\input{../tex-inputs/latex-leseansicht-abspann}
